\documentclass[a4paper,10pt]{letter}

\usepackage[T1]{fontenc}
\usepackage[utf8]{inputenc}
\usepackage{lmodern}
\usepackage[hidelinks]{hyperref}
\usepackage{microtype}
\usepackage[margin=1in]{geometry}

\signature{Omar Iskandarani\\
Independent Researcher, Groningen, The Netherlands\\
ORCID: 0009-0006-1686-3961\\
Email: \href{mailto:info@omariskandarani.com}{info@omariskandarani.com}}

\address{Omar Iskandarani\\
Hopstraat 11\\
9903 BV Appingedaam\\
The Netherlands}

\date{\today}

\begin{document}
    
    \begin{letter}{Editors\\ \textit{International Journal of Theoretical Physics}}

\opening{Dear Editor,}

Please consider the enclosed manuscript,
\textit{Relational Time-of-Arrival as a Covariant Field Observable: Integrable Event Currents, Clock Holonomy, and a Continuum Clock Limit},
for publication as a Research Article in \textit{International Journal of Theoretical Physics}.

The manuscript develops a covariant theoretical framework for quantum time-of-arrival observables. Instead of introducing a universal self-adjoint time operator, it formulates time-of-arrival as a detector-level field observable built from two conserved structures: a matter current $J^\mu$ describing detector flux through a world-tube $\Sigma$, and an effective event-count current $j^\mu_{\mathrm{ev}}$ defining a physical clock sector.

The main structural contribution is the separation of current conservation from local integrability. Conservation stabilizes event counts, whereas local integrability of the coarse-grained event-current one-form is required for a scalar continuum clock field to exist. The obstruction is the event-current vorticity $\Omega^{\mathrm{ev}}_{\mu\nu}$, which leads in the non-integrable sector to path-conditioned time-of-arrival distributions and clock-holonomy contributions.

The paper also provides a worked $(1+1)$-dimensional Klein--Gordon example with a gapped quadratic clock effective field theory. In the integrable spectator-clock regime, the time-of-arrival density reduces to a convolution of the semiclassical flux profile with the clock readout distribution. For Gaussian detector smearing, the regulated clock variance is evaluated explicitly as
\[
    \sigma_{\tau,\ell}^{2} = (4\pi)^{-1} K_0(\mu_\tau^2 \ell^2),
\]
yielding a parameter-controlled contribution to arrival-time broadening.

I believe the manuscript is suitable for \textit{International Journal of Theoretical Physics} because it addresses a theoretical problem at the intersection of quantum measurement theory, quantum field theory, relational time, and covariant observables. The work is self-contained, does not rely on speculative physical assumptions, and is positioned within established literature on time-of-arrival, relational clocks, partial observables, and field-theoretic detector descriptions.

The submission is original, has not been published previously, and is not under consideration elsewhere. A preprint version is available on Zenodo:
\href{https://doi.org/10.5281/zenodo.20486634}{10.5281/zenodo.20486634}.
No datasets were generated or analyzed during the current study. There is no external funding. On behalf of all authors, the corresponding author states that there is no conflict of interest.

Thank you for your consideration.

\closing{Sincerely,}
    
    \end{letter}

\end{document}